\section{V}

% verändern (to change)
\begin{verbentry}{
  style = {acc},
  toc = {verändern (to change)},
  name = {verändern},
  english = {to change},
  type = {irregular, + Akkusativ},
  auxiliary = {haben}
}
 
    
     \vspace{5pt}

    \hrule
    
    \vspace{5pt}

  \begin{tabular}{rl}
     \multicolumn{2}{l}{\textbf{PRÄSENS} }\\
    ich & \textbf{verändere}\\
    du & \textbf{veränderst}\\
    er/sie/es & \textbf{verändert}\\
    wir & \textbf{verändern}\\
    ihr & \textbf{verändert}\\
    sie/Sie & \textbf{verändern}\\
  \end{tabular}
  \hfill
     \begin{tabular}{rl}
    \multicolumn{2}{l}{\textbf{PRÄTERITUM} }\\
    ich & \textbf{veränderte}\\
    du & \textbf{verändertest}\\
    er/sie/es & \textbf{veränderte}\\
    wir & \textbf{veränderten}\\
    ihr & \textbf{verändertet}\\
    sie/Sie & \textbf{veränderten}\\
  \end{tabular}
  \hfill
  \begin{tabular}{rl}
     \multicolumn{2}{l}{\textbf{IMPERATIV} }\\
   & \textbf{veränder(e) (du)!}\\
   & \textbf{verändern wir!}\\
   & \textbf{verändert (ihr)!}\\
   & \textbf{verändern Sie!}\\
   & \vspace{10pt}
  \end{tabular}

    \vspace{10pt}

 \begin{tabular}{rl}
     \multicolumn{2}{l}{\textbf{PERFEKT}}\\
  &  haben + \textbf{verändert}\\
  \end{tabular}
\hfill
   \begin{tabular}{rl}
   
    \multicolumn{2}{l}{\textbf{FUTURE I} }\\
  &  werden + \textbf{verändern}\\
 
  \end{tabular}
\hfill
   \begin{tabular}{rl}
      \multicolumn{2}{l}{\textbf{PLUSQUAMPERFEKT}}\\
   & hatten + \textbf{verändert}\\
  \end{tabular}
  
  
\vspace{10pt}

\begin{tabular}{lll}
\multicolumn{3}{l}{\textbf{MIT PRÄPOSITIONEN}}\\
& \textbf{---} & \\
\end{tabular}
\hfill
\begin{tabular}{lll}
\multicolumn{3}{l}{\textbf{TRENNBARE VERBEN}}\\
& \textbf{---} & \\
\end{tabular}
\vspace{-5pt}
\end{verbentry}

% verbinden (to connect / link)
\begin{verbentry}{
  style = {default},
  toc = {verbinden (to connect / link)},
  name = {verbinden},
  english = {to connect / link},
  type = {regular},
  auxiliary = {haben}
}


\vspace{5pt}
\hrule
\vspace{5pt}

\begin{tabular}{rl}
\multicolumn{2}{l}{\textbf{PRÄSENS}}\\
ich & \textbf{verbinde}\\
du & \textbf{verbindest}\\
er/sie/es & \textbf{verbindet}\\
wir & \textbf{verbinden}\\
ihr & \textbf{verbindet}\\
sie/Sie & \textbf{verbinden}\\
\end{tabular}
\hfill
\begin{tabular}{rl}
\multicolumn{2}{l}{\textbf{PRÄTERITUM}}\\
ich & \textbf{verband}\\
du & \textbf{verbandst}\\
er/sie/es & \textbf{verband}\\
wir & \textbf{verbanden}\\
ihr & \textbf{verbandet}\\
sie/Sie & \textbf{verbanden}\\
\end{tabular}
\hfill
\begin{tabular}{rl}
\multicolumn{2}{l}{\textbf{IMPERATIV}}\\
& \textbf{verbinde (du)!}\\
& \textbf{verbinden wir!}\\
& \textbf{verbindet (ihr)!}\\
& \textbf{verbinden Sie!}\\
\end{tabular}

\vspace{10pt}

\begin{tabular}{rl}
\multicolumn{2}{l}{\textbf{PERFEKT}}\\
& haben + \textbf{verbunden}\\
\end{tabular}
\hfill
\begin{tabular}{rl}
\multicolumn{2}{l}{\textbf{FUTURE I}}\\
& werden + \textbf{verbinden}\\
\end{tabular}
\hfill
\begin{tabular}{rl}
\multicolumn{2}{l}{\textbf{PLUSQUAMPERFEKT}}\\
& hatten + \textbf{verbunden}\\
\end{tabular}

\vspace{10pt}

\begin{tabular}{lll}
\multicolumn{3}{l}{\textbf{MIT PRÄPOSITIONEN}}\\
& \textbf{verbinden mit} + Dat & (to connect with)\\
& \textbf{verbinden zu} + Dat & (to link to / form into)\\
\end{tabular}
\hfill
\begin{tabular}{lll}
\multicolumn{3}{l}{\textbf{TRENNBARE VERBEN}}\\
\vspace{20pt}
\end{tabular}
\vspace{-5pt}
\end{verbentry}

% verbreiten (to spread / distribute)
\begin{verbentry}{
  style = {default},
  toc = {verbreiten (to spread / distribute)},
  name = {verbreiten},
  english = {to spread / distribute},
  type = {regular},
  auxiliary = {haben}
}


\vspace{5pt}
\hrule
\vspace{5pt}

\begin{tabular}{rl}
\multicolumn{2}{l}{\textbf{PRÄSENS}}\\
ich & \textbf{verbreite}\\
du & \textbf{verbreitest}\\
er/sie/es & \textbf{verbreitet}\\
wir & \textbf{verbreiten}\\
ihr & \textbf{verbreitet}\\
sie/Sie & \textbf{verbreiten}\\
\end{tabular}
\hfill
\begin{tabular}{rl}
\multicolumn{2}{l}{\textbf{PRÄTERITUM}}\\
ich & \textbf{verbreitete}\\
du & \textbf{verbreitetest}\\
er/sie/es & \textbf{verbreitete}\\
wir & \textbf{verbreiteten}\\
ihr & \textbf{verbreitetet}\\
sie/Sie & \textbf{verbreiteten}\\
\end{tabular}
\hfill
\begin{tabular}{rl}
\multicolumn{2}{l}{\textbf{IMPERATIV}}\\
& \textbf{verbreite (du)!}\\
& \textbf{verbreiten wir!}\\
& \textbf{verbreitet (ihr)!}\\
& \textbf{verbreiten Sie!}\\
\end{tabular}

\vspace{10pt}

\begin{tabular}{rl}
\multicolumn{2}{l}{\textbf{PERFEKT}}\\
& haben + \textbf{verbreitet}\\
\end{tabular}
\hfill
\begin{tabular}{rl}
\multicolumn{2}{l}{\textbf{FUTURE I}}\\
& werden + \textbf{verbreiten}\\
\end{tabular}
\hfill
\begin{tabular}{rl}
\multicolumn{2}{l}{\textbf{PLUSQUAMPERFEKT}}\\
& hatten + \textbf{verbreitet}\\
\end{tabular}

\vspace{10pt}

\begin{tabular}{lll}
\multicolumn{3}{l}{\textbf{MIT PRÄPOSITIONEN}}\\
& \textbf{verbreiten unter} + Dat & (to spread among)\\
& \textbf{verbreiten über} + Akk & (to spread about)\\
& \textbf{verbreiten durch} + Akk & (to spread through)\\
\end{tabular}
\hfill
\begin{tabular}{lll}
\multicolumn{3}{l}{\textbf{TRENNBARE VERBEN}}\\
\vspace{30pt}\end{tabular}
\vspace{-5pt}
\end{verbentry}

% verbringen (to spend time)
\begin{verbentry}{
  style = {acc},
  toc = {verbringen (to spend time)},
  name = {verbringen},
  english = {to spend time},
  type = {irregular, + Akkusativ},
  auxiliary = {haben}
}


\vspace{5pt}
\hrule
\vspace{5pt}

\begin{tabular}{rl}
\multicolumn{2}{l}{\textbf{PRÄSENS}}\\
ich & \textbf{verbringe}\\
du & \textbf{verbringst}\\
er/sie/es & \textbf{verbringt}\\
wir & \textbf{verbringen}\\
ihr & \textbf{verbringt}\\
sie/Sie & \textbf{verbringen}\\
\end{tabular}
\hfill
\begin{tabular}{rl}
\multicolumn{2}{l}{\textbf{PRÄTERITUM}}\\
ich & \textbf{verbrachte}\\
du & \textbf{verbrachtest}\\
er/sie/es & \textbf{verbrachte}\\
wir & \textbf{verbrachten}\\
ihr & \textbf{verbrachtet}\\
sie/Sie & \textbf{verbrachten}\\
\end{tabular}
\hfill
\begin{tabular}{rl}
\multicolumn{2}{l}{\textbf{IMPERATIV}}\\
& \textbf{verbring(e) (du)!}\\
& \textbf{verbringen wir!}\\
& \textbf{verbringt (ihr)!}\\
& \textbf{verbringen Sie!}\\
\end{tabular}

\vspace{10pt}

\begin{tabular}{rl}
\multicolumn{2}{l}{\textbf{PERFEKT}}\\
& haben + \textbf{verbracht}\\
\end{tabular}
\hfill
\begin{tabular}{rl}
\multicolumn{2}{l}{\textbf{FUTURE I}}\\
& werden + \textbf{verbringen}\\
\end{tabular}
\hfill
\begin{tabular}{rl}
\multicolumn{2}{l}{\textbf{PLUSQUAMPERFEKT}}\\
& hatten + \textbf{verbracht}\\
\end{tabular}

\vspace{10pt}

\begin{tabular}{lll}
\multicolumn{3}{l}{\textbf{MIT PRÄPOSITIONEN}}\\
& \textbf{---} & \\
\end{tabular}
\hfill
\begin{tabular}{lll}
\multicolumn{3}{l}{\textbf{TRENNBARE VERBEN}}\\
& \textbf{---} & \\
\end{tabular}

\vspace{-5pt}
\end{verbentry}

% vereinbaren (to arrange / agree on)
\begin{verbentry}{
  style = {default},
  toc = {vereinbaren (to arrange / agree on)},
  name = {vereinbaren},
  english = {to arrange / agree on},
  type = {regular},
  auxiliary = {haben}
}


\vspace{5pt}
\hrule
\vspace{5pt}

\begin{tabular}{rl}
\multicolumn{2}{l}{\textbf{PRÄSENS}}\\
ich & \textbf{vereinbare}\\
du & \textbf{vereinbarst}\\
er/sie/es & \textbf{vereinbart}\\
wir & \textbf{vereinbaren}\\
ihr & \textbf{vereinbart}\\
sie/Sie & \textbf{vereinbaren}\\
\end{tabular}
\hfill
\begin{tabular}{rl}
\multicolumn{2}{l}{\textbf{PRÄTERITUM}}\\
ich & \textbf{vereinbarte}\\
du & \textbf{vereinbartest}\\
er/sie/es & \textbf{vereinbarte}\\
wir & \textbf{vereinbarten}\\
ihr & \textbf{vereinbartet}\\
sie/Sie & \textbf{vereinbarten}\\
\end{tabular}
\hfill
\begin{tabular}{rl}
\multicolumn{2}{l}{\textbf{IMPERATIV}}\\
& \textbf{vereinbare (du)!}\\
& \textbf{vereinbaren wir!}\\
& \textbf{vereinbart (ihr)!}\\
& \textbf{vereinbaren Sie!}\\
\end{tabular}

\vspace{10pt}

\begin{tabular}{rl}
\multicolumn{2}{l}{\textbf{PERFEKT}}\\
& haben + \textbf{vereinbart}\\
\end{tabular}
\hfill
\begin{tabular}{rl}
\multicolumn{2}{l}{\textbf{FUTURE I}}\\
& werden + \textbf{vereinbaren}\\
\end{tabular}
\hfill
\begin{tabular}{rl}
\multicolumn{2}{l}{\textbf{PLUSQUAMPERFEKT}}\\
& hatten + \textbf{vereinbart}\\
\end{tabular}

\vspace{10pt}

\begin{tabular}{lll}
\multicolumn{3}{l}{\textbf{MIT PRÄPOSITIONEN}}\\
& \textbf{vereinbaren mit} + Dat & (to arrange with / agree with someone)\\
& \textbf{vereinbaren über} + Akk & (to agree on / about something)\\
\end{tabular}
\hfill
\begin{tabular}{lll}
\multicolumn{3}{l}{\textbf{TRENNBARE VERBEN}}\\
\vspace{20pt}\end{tabular}
\vspace{-5pt}
\end{verbentry}

% vergleichen (to compare)
\begin{verbentry}{
  style = {dat},
  toc = {vergleichen (to compare)},
  name = {vergleichen},
  english = {to compare},
  type = {irregular, + Accuative + Dativ},
  auxiliary = {haben}
}
 
    
     \vspace{5pt}

    \hrule
    
    \vspace{5pt}

  \begin{tabular}{rl}
     \multicolumn{2}{l}{\textbf{PRÄSENS} }\\
    ich & \textbf{vergleiche}\\
    du & \textbf{vergleichst}\\
    er/sie/es & \textbf{vergleicht}\\
    wir & \textbf{vergleichen}\\
    ihr & \textbf{vergleicht}\\
    sie/Sie & \textbf{vergleichen}\\
  \end{tabular}
  \hfill
     \begin{tabular}{rl}
    \multicolumn{2}{l}{\textbf{PRÄTERITUM} }\\
    ich & \textbf{verglich}\\
    du & \textbf{verglichst}\\
    er/sie/es & \textbf{verglich}\\
    wir & \textbf{verglichen}\\
    ihr & \textbf{verglicht}\\
    sie/Sie & \textbf{verglichen}\\
  \end{tabular}
  \hfill
  \begin{tabular}{rl}
     \multicolumn{2}{l}{\textbf{IMPERATIV} }\\
   & \textbf{vergleich(e) (du)!}\\
   & \textbf{vergleichen wir!}\\
   & \textbf{vergleicht (ihr)!}\\
   & \textbf{vergleichen Sie!}\\
   & \vspace{10pt}
  \end{tabular}

    \vspace{10pt}

 \begin{tabular}{rl}
     \multicolumn{2}{l}{\textbf{PERFEKT}}\\
  &  haben + \textbf{verglichen}\\
  \end{tabular}
\hfill
   \begin{tabular}{rl}
   
    \multicolumn{2}{l}{\textbf{FUTURE I} }\\
  &  werden + \textbf{vergleichen}\\
 
  \end{tabular}
\hfill
   \begin{tabular}{rl}
      \multicolumn{2}{l}{\textbf{PLUSQUAMPERFEKT}}\\
  &  hatten + \textbf{verglichen}\\
  \end{tabular}
  
  
\vspace{10pt}

\begin{tabular}{lll}
\multicolumn{3}{l}{\textbf{MIT PRÄPOSITIONEN}}\\
& \textbf{---} & \\
\end{tabular}
\hfill
\begin{tabular}{lll}
\multicolumn{3}{l}{\textbf{TRENNBARE VERBEN}}\\
& \textbf{---} & \\
\end{tabular}
\vspace{-5pt}
\end{verbentry}

% vergessen (to forget)
\begin{verbentry}{
  style = {acc},
  toc = {vergessen (to forget)},
  name = {vergessen},
  english = {to forget},
  type = {irregular, + Akkusativ},
  auxiliary = {haben}
}
 
    
     \vspace{5pt}

    \hrule
    
    \vspace{5pt}

  \begin{tabular}{rl}
     \multicolumn{2}{l}{\textbf{PRÄSENS} }\\
    ich & \textbf{vergesse}\\
    du & \textbf{vergisst}\\
    er/sie/es & \textbf{vergisst}\\
    wir & \textbf{vergessen}\\
    ihr & \textbf{vergesst}\\
    sie/Sie & \textbf{vergessen}\\
  \end{tabular}
  \hfill
     \begin{tabular}{rl}
    \multicolumn{2}{l}{\textbf{PRÄTERITUM} }\\
    ich & \textbf{vergaß}\\
    du & \textbf{vergaßest}\\
    er/sie/es & \textbf{vergaß}\\
    wir & \textbf{vergaßen}\\
    ihr & \textbf{vergaßt}\\
    sie/Sie & \textbf{vergaßen}\\
  \end{tabular}
  \hfill
  \begin{tabular}{rl}
     \multicolumn{2}{l}{\textbf{IMPERATIV} }\\
   & \textbf{vergiss (du)!}\\
   & \textbf{vergessen wir!}\\
   & \textbf{vergesst (ihr)!}\\
   & \textbf{vergessen Sie!}\\
   & \vspace{10pt}
  \end{tabular}

    \vspace{10pt}

 \begin{tabular}{rl}
     \multicolumn{2}{l}{\textbf{PERFEKT}}\\
  &  haben + \textbf{vergessen}\\
  \end{tabular}
\hfill
   \begin{tabular}{rl}
   
    \multicolumn{2}{l}{\textbf{FUTURE I} }\\
  &  werden + \textbf{vergessen}\\
 
  \end{tabular}
\hfill
   \begin{tabular}{rl}
      \multicolumn{2}{l}{\textbf{PLUSQUAMPERFEKT}}\\
  &  hatten + \textbf{vergessen}\\
  \end{tabular}
  
  
\vspace{10pt}

\begin{tabular}{lll}
\multicolumn{3}{l}{\textbf{MIT PRÄPOSITIONEN}}\\
& \textbf{---} & \\
\end{tabular}
\hfill
\begin{tabular}{lll}
\multicolumn{3}{l}{\textbf{TRENNBARE VERBEN}}\\
& \textbf{---} & \\
\end{tabular}
\vspace{-5pt}
\end{verbentry}

% verhindern (to prevent)
\begin{verbentry}{
  style = {default},
  toc = {verhindern (to prevent)},
  name = {verhindern},
  english = {to prevent},
  type = {regular, inseparable},
  auxiliary = {haben}
}
 
    
     \vspace{5pt}

    \hrule
    
    \vspace{5pt}

  \begin{tabular}{rl}
     \multicolumn{2}{l}{\textbf{PRÄSENS} }\\
    ich & \textbf{verhindere}\\
    du & \textbf{verhinderst}\\
    er/sie/es & \textbf{verhindert}\\
    wir & \textbf{verhindern}\\
    ihr & \textbf{verhindert}\\
    sie/Sie & \textbf{verhindern}\\
  \end{tabular}
  \hfill
     \begin{tabular}{rl}
    \multicolumn{2}{l}{\textbf{PRÄTERITUM} }\\
    ich & \textbf{verhinderte}\\
    du & \textbf{verhindertest}\\
    er/sie/es & \textbf{verhinderte}\\
    wir & \textbf{verhinderten}\\
    ihr & \textbf{verhindertet}\\
    sie/Sie & \textbf{verhinderten}\\
  \end{tabular}
  \hfill
  \begin{tabular}{rl}
     \multicolumn{2}{l}{\textbf{IMPERATIV} }\\
   & \textbf{verhindere (du)!}\\
   & \textbf{verhindern wir!}\\
   & \textbf{verhindert (ihr)!}\\
   & \textbf{verhindern Sie!}\\
   & \vspace{10pt}
  \end{tabular}

    \vspace{10pt}

 \begin{tabular}{rl}
     \multicolumn{2}{l}{\textbf{PERFEKT}}\\
   & haben + \textbf{verhindert}\\
  \end{tabular}
\hfill
   \begin{tabular}{rl}
    \multicolumn{2}{l}{\textbf{FUTURE I} }\\
   & werden + \textbf{verhindern}\\
  \end{tabular}
\hfill
   \begin{tabular}{rl}
      \multicolumn{2}{l}{\textbf{PLUSQUAMPERFEKT}}\\
   & hatten + \textbf{verhindert}\\
  \end{tabular}

   \vspace{10pt}

   \begin{tabular}{lll}
      \multicolumn{3}{l}{\textbf{MIT PRÄPOSITIONEN}}\\
& \textbf{verhindern an} + Dat & (to prevent from)\\
& \textbf{verhindern durch} + Akk & (to prevent by means of)\\
\end{tabular}
  \hfill
   \begin{tabular}{lll}
\multicolumn{3}{l}{\textbf{TRENNBARE VERBEN}}\\
& \textbf{---} & \\
\end{tabular}

  \vspace{-5pt}
\end{verbentry}

% verleihen (to lend / award)
\begin{verbentry}{
  style = {default},
  toc = {verleihen (to lend / award)},
  name = {verleihen},
  english = {to lend / award},
  type = {regular},
  auxiliary = {haben}
}
 
    
     \vspace{5pt}

    \hrule
    
    \vspace{5pt}

  \begin{tabular}{rl}
     \multicolumn{2}{l}{\textbf{PRÄSENS} }\\
    ich & \textbf{verleihe}\\
    du & \textbf{verleihst}\\
    er/sie/es & \textbf{verleiht}\\
    wir & \textbf{verleihen}\\
    ihr & \textbf{verleiht}\\
    sie/Sie & \textbf{verleihen}\\
  \end{tabular}
  \hfill
  \begin{tabular}{rl}
    \multicolumn{2}{l}{\textbf{PRÄTERITUM} }\\
    ich & \textbf{verlieh}\\
    du & \textbf{verliehst}\\
    er/sie/es & \textbf{verlieh}\\
    wir & \textbf{verliehen}\\
    ihr & \textbf{verlieht}\\
    sie/Sie & \textbf{verliehen}\\
  \end{tabular}
  \hfill
  \begin{tabular}{rl}
     \multicolumn{2}{l}{\textbf{IMPERATIV} }\\
   & \textbf{verleihe (du)!}\\
   & \textbf{verleihen wir!}\\
   & \textbf{verleiht (ihr)!}\\
   & \textbf{verleihen Sie!}\\
   & \vspace{10pt}
  \end{tabular}

 \vspace{10pt}

 \begin{tabular}{rl}
     \multicolumn{2}{l}{\textbf{PERFEKT}}\\
   & haben + \textbf{verliehen}\\
  \end{tabular}
\hfill
\begin{tabular}{rl}
    \multicolumn{2}{l}{\textbf{FUTURE I}}\\
   & werden + \textbf{verleihen}\\
  \end{tabular}
\hfill
\begin{tabular}{rl}
    \multicolumn{2}{l}{\textbf{PLUSQUAMPERFEKT}}\\
   & hatten + \textbf{verliehen}\\
  \end{tabular}

 \vspace{10pt}

\begin{tabular}{lll}
      \multicolumn{3}{l}{\textbf{MIT PRÄPOSITIONEN}}\\
& \textbf{verleihen an} + Akk & (to lend to)\\
& \textbf{verleihen für} + Akk & (to award for)\\
\end{tabular}
  \hfill
\begin{tabular}{lll}
      \multicolumn{3}{l}{\textbf{TRENNBARE VERBEN}}\\
 & \textbf{ver$|$leihen} & (to lend out / award)\\
\vspace{10pt}
  \end{tabular}


\vspace{-5pt}
\end{verbentry}

% verletzen (to injure / to hurt)
\begin{verbentry}{
  style = {default},
  toc = {verletzen (to injure / to hurt)},
  name = {verletzen},
  english = {to injure / to hurt},
  type = {regular},
  auxiliary = {haben}
}
 
    
     \vspace{5pt}

    \hrule
    
    \vspace{5pt}

  \begin{tabular}{rl}
     \multicolumn{2}{l}{\textbf{PRÄSENS} }\\
    ich & \textbf{verletze}\\
    du & \textbf{verletzt}\\
    er/sie/es & \textbf{verletzt}\\
    wir & \textbf{verletzen}\\
    ihr & \textbf{verletzt}\\
    sie/Sie & \textbf{verletzen}\\
  \end{tabular}
  \hfill
     \begin{tabular}{rl}
    \multicolumn{2}{l}{\textbf{PRÄTERITUM} }\\
    ich & \textbf{verletzte}\\
    du & \textbf{verletztest}\\
    er/sie/es & \textbf{verletzte}\\
    wir & \textbf{verletzten}\\
    ihr & \textbf{verletztet}\\
    sie/Sie & \textbf{verletzten}\\
  \end{tabular}
  \hfill
  \begin{tabular}{rl}
     \multicolumn{2}{l}{\textbf{IMPERATIV} }\\
   & \textbf{verletz(e) (du)!}\\
   & \textbf{verletzen wir!}\\
   & \textbf{verletzt (ihr)!}\\
   & \textbf{verletzen Sie!}\\
   & \vspace{10pt}
  \end{tabular}

    \vspace{10pt}

 \begin{tabular}{rl}
     \multicolumn{2}{l}{\textbf{PERFEKT}}\\
   & haben + \textbf{verletzt}\\
  \end{tabular}
\hfill
   \begin{tabular}{rl}
    \multicolumn{2}{l}{\textbf{FUTURE I} }\\
   & werde + \textbf{verletzen}\\
  \end{tabular}
\hfill
   \begin{tabular}{rl}
      \multicolumn{2}{l}{\textbf{PLUSQUAMPERFEKT}}\\
   & hatten + \textbf{verletzt}\\
  \end{tabular}

   \vspace{10pt}

   \begin{tabular}{lll}
      \multicolumn{3}{l}{\textbf{MIT PRÄPOSITIONEN}}\\
& \textbf{sich verletzen an} + Dat & (to injure oneself on)\\
& \textbf{verletzen mit} + Dat & (to injure with)\\
\end{tabular}
  \hfill
   \begin{tabular}{lll}
\multicolumn{3}{l}{\textbf{TRENNBARE VERBEN}}\\
& \textbf{---} & \\
\end{tabular}

  \vspace{-5pt}
\end{verbentry}

% verlieren (to lose)
\begin{verbentry}{
  style = {default},
  toc = {verlieren (to lose)},
  name = {verlieren},
  english = {to lose},
  type = {irregular},
  auxiliary = {haben}
}
 
    
     \vspace{5pt}

    \hrule
    
    \vspace{5pt}

  \begin{tabular}{rl}
     \multicolumn{2}{l}{\textbf{PRÄSENS} }\\
    ich & \textbf{verliere}\\
    du & \textbf{verlierst}\\
    er/sie/es & \textbf{verliert}\\
    wir & \textbf{verlieren}\\
    ihr & \textbf{verliert}\\
    sie/Sie & \textbf{verlieren}\\
  \end{tabular}
  \hfill
     \begin{tabular}{rl}
    \multicolumn{2}{l}{\textbf{PRÄTERITUM} }\\
    ich & \textbf{verlor}\\
    du & \textbf{verlorst}\\
    er/sie/es & \textbf{verlor}\\
    wir & \textbf{verloren}\\
    ihr & \textbf{verlort}\\
    sie/Sie & \textbf{verloren}\\
  \end{tabular}
  \hfill
  \begin{tabular}{rl}
     \multicolumn{2}{l}{\textbf{IMPERATIV} }\\
   & \textbf{verlier(e) (du)!}\\
   & \textbf{verlieren wir!}\\
   & \textbf{verliert (ihr)!}\\
   & \textbf{verlieren Sie!}\\
   & \vspace{10pt}
  \end{tabular}

    \vspace{10pt}

 \begin{tabular}{rl}
     \multicolumn{2}{l}{\textbf{PERFEKT}}\\
   & haben + \textbf{verloren}\\
  \end{tabular}
\hfill
   \begin{tabular}{rl}
    \multicolumn{2}{l}{\textbf{FUTURE I} }\\
   & werde + \textbf{verlieren}\\
  \end{tabular}
\hfill
   \begin{tabular}{rl}
      \multicolumn{2}{l}{\textbf{PLUSQUAMPERFEKT}}\\
   & hatten + \textbf{verloren}\\
  \end{tabular}

   \vspace{10pt}

   \begin{tabular}{lll}
      \multicolumn{3}{l}{\textbf{MIT PRÄPOSITIONEN}}\\
& \textbf{verlieren an} + Akk & (to lose to)\\
& \textbf{verlieren durch} + Akk & (to lose through / because of)\\
\end{tabular}
  \hfill
   \begin{tabular}{lll}
\multicolumn{3}{l}{\textbf{TRENNBARE VERBEN}}\\
& \textbf{---} & \\
\end{tabular}

  \vspace{-5pt}
\end{verbentry}

% vermeiden (to avoid)
\begin{verbentry}{
  style = {default},
  toc = {vermeiden (to avoid)},
  name = {vermeiden},
  english = {to avoid},
  type = {irregular},
  auxiliary = {haben}
}


\vspace{5pt}
\hrule
\vspace{5pt}

\begin{tabular}{rl}
\multicolumn{2}{l}{\textbf{PRÄSENS}}\\
ich & \textbf{vermeide}\\
du & \textbf{vermeidest}\\
er/sie/es & \textbf{vermeidet}\\
wir & \textbf{vermeiden}\\
ihr & \textbf{vermeidet}\\
sie/Sie & \textbf{vermeiden}\\
\end{tabular}
\hfill
\begin{tabular}{rl}
\multicolumn{2}{l}{\textbf{PRÄTERITUM}}\\
ich & \textbf{vermied}\\
du & \textbf{vermiedest}\\
er/sie/es & \textbf{vermied}\\
wir & \textbf{vermieden}\\
ihr & \textbf{vermiedet}\\
sie/Sie & \textbf{vermieden}\\
\end{tabular}
\hfill
\begin{tabular}{rl}
\multicolumn{2}{l}{\textbf{IMPERATIV}}\\
& \textbf{vermeide (du)!}\\
& \textbf{vermeiden wir!}\\
& \textbf{vermeidet (ihr)!}\\
& \textbf{vermeiden Sie!}\\
\end{tabular}

\vspace{10pt}

\begin{tabular}{rl}
\multicolumn{2}{l}{\textbf{PERFEKT}}\\
& haben + \textbf{vermieden}\\
\end{tabular}
\hfill
\begin{tabular}{rl}
\multicolumn{2}{l}{\textbf{FUTURE I}}\\
& werden + \textbf{vermeiden}\\
\end{tabular}
\hfill
\begin{tabular}{rl}
\multicolumn{2}{l}{\textbf{PLUSQUAMPERFEKT}}\\
& hatten + \textbf{vermieden}\\
\end{tabular}

\vspace{10pt}

\begin{tabular}{lll}
\multicolumn{3}{l}{\textbf{MIT PRÄPOSITIONEN}}\\
& \textbf{---} & \\
\end{tabular}
\hfill
\begin{tabular}{lll}
\multicolumn{3}{l}{\textbf{TRENNBARE VERBEN}}\\
& \textbf{---} & \\
\end{tabular}
\vspace{-5pt}
\end{verbentry}

% vermissen (to miss someone)
\begin{verbentry}{
  style = {default},
  toc = {vermissen (to miss someone)},
  name = {vermissen},
  english = {to miss someone},
  type = {regular},
  auxiliary = {haben}
}


\vspace{5pt}
\hrule
\vspace{5pt}

\begin{tabular}{rl}
\multicolumn{2}{l}{\textbf{PRÄSENS}}\\
ich & \textbf{vermisse}\\
du & \textbf{vermisst}\\
er/sie/es & \textbf{vermisst}\\
wir & \textbf{vermissen}\\
ihr & \textbf{vermisst}\\
sie/Sie & \textbf{vermissen}\\
\end{tabular}
\hfill
\begin{tabular}{rl}
\multicolumn{2}{l}{\textbf{PRÄTERITUM}}\\
ich & \textbf{vermisste}\\
du & \textbf{vermisstest}\\
er/sie/es & \textbf{vermisste}\\
wir & \textbf{vermissten}\\
ihr & \textbf{vermisstet}\\
sie/Sie & \textbf{vermissten}\\
\end{tabular}
\hfill
\begin{tabular}{rl}
\multicolumn{2}{l}{\textbf{IMPERATIV}}\\
& \textbf{vermisse (du)!}\\
& \textbf{vermissen wir!}\\
& \textbf{vermisst (ihr)!}\\
& \textbf{vermissen Sie!}\\
\end{tabular}

\vspace{10pt}

\begin{tabular}{rl}
\multicolumn{2}{l}{\textbf{PERFEKT}}\\
& haben + \textbf{vermisst}\\
\end{tabular}
\hfill
\begin{tabular}{rl}
\multicolumn{2}{l}{\textbf{FUTURE I}}\\
& werden + \textbf{vermissen}\\
\end{tabular}
\hfill
\begin{tabular}{rl}
\multicolumn{2}{l}{\textbf{PLUSQUAMPERFEKT}}\\
& hatten + \textbf{vermisst}\\
\end{tabular}

\vspace{10pt}

\begin{tabular}{lll}
\multicolumn{3}{l}{\textbf{MIT PRÄPOSITIONEN}}\\
& \textbf{vermissen} + Akk & (to miss someone/something)\\
\end{tabular}
\hfill
\begin{tabular}{lll}
\multicolumn{3}{l}{\textbf{TRENNBARE VERBEN}}\\
& \textbf{---} & \\
\end{tabular}
\vspace{-5pt}
\end{verbentry}

% verputzen (to plaster / render)
\begin{verbentry}{
  style = {default},
  toc = {verputzen (to plaster / render)},
  name = {verputzen},
  english = {to plaster / render},
  type = {regular},
  auxiliary = {haben}
}
 
    
     \vspace{5pt}

    \hrule
    
    \vspace{5pt}

  \begin{tabular}{rl}
     \multicolumn{2}{l}{\textbf{PRÄSENS} }\\
    ich & \textbf{verputze}\\
    du & \textbf{verputzt}\\
    er/sie/es & \textbf{verputzt}\\
    wir & \textbf{verputzen}\\
    ihr & \textbf{verputzt}\\
    sie/Sie & \textbf{verputzen}\\
  \end{tabular}
  \hfill
  \begin{tabular}{rl}
    \multicolumn{2}{l}{\textbf{PRÄTERITUM} }\\
    ich & \textbf{verputzte}\\
    du & \textbf{verputztest}\\
    er/sie/es & \textbf{verputzte}\\
    wir & \textbf{verputzten}\\
    ihr & \textbf{verputztet}\\
    sie/Sie & \textbf{verputzten}\\
  \end{tabular}
  \hfill
  \begin{tabular}{rl}
     \multicolumn{2}{l}{\textbf{IMPERATIV} }\\
   & \textbf{verputze (du)!}\\
   & \textbf{verputzen wir!}\\
   & \textbf{verputzt (ihr)!}\\
   & \textbf{verputzen Sie!}\\
   & \vspace{10pt}
  \end{tabular}

 \vspace{10pt}

 \begin{tabular}{rl}
     \multicolumn{2}{l}{\textbf{PERFEKT}}\\
   & haben + \textbf{verputzt}\\
  \end{tabular}
\hfill
\begin{tabular}{rl}
    \multicolumn{2}{l}{\textbf{FUTURE I}}\\
   & werden + \textbf{verputzen}\\
  \end{tabular}
\hfill
\begin{tabular}{rl}
    \multicolumn{2}{l}{\textbf{PLUSQUAMPERFEKT}}\\
   & hatten + \textbf{verputzt}\\
  \end{tabular}

 \vspace{10pt}

\begin{tabular}{lll}
      \multicolumn{3}{l}{\textbf{MIT PRÄPOSITIONEN}}\\
& \textbf{verputzen mit} + Dat & (to plaster with)\\
\end{tabular}
  \hfill
\begin{tabular}{lll}
      \multicolumn{3}{l}{\textbf{TRENNBARE VERBEN}}\\
 & \textbf{ein$|$verputzen} & (to plaster in / apply)\\
  \end{tabular}


\vspace{-5pt}
\end{verbentry}

% verschreiben (to prescribe)
\begin{verbentry}{
  style = {default},
  toc = {verschreiben (to prescribe)},
  name = {verschreiben},
  english = {to prescribe},
  type = {irregular},
  auxiliary = {haben / sich + haben}
}


\vspace{5pt}
\hrule
\vspace{5pt}

\begin{tabular}{rl}
\multicolumn{2}{l}{\textbf{PRÄSENS}}\\
ich & \textbf{verschreibe}\\
du & \textbf{verschreibst}\\
er/sie/es & \textbf{verschreibt}\\
wir & \textbf{verschreiben}\\
ihr & \textbf{verschreibt}\\
sie/Sie & \textbf{verschreiben}\\
\end{tabular}
\hfill
\begin{tabular}{rl}
\multicolumn{2}{l}{\textbf{PRÄTERITUM}}\\
ich & \textbf{verschrieb}\\
du & \textbf{verschriebst}\\
er/sie/es & \textbf{verschrieb}\\
wir & \textbf{verschrieben}\\
ihr & \textbf{verschriebt}\\
sie/Sie & \textbf{verschrieben}\\
\end{tabular}
\hfill
\begin{tabular}{rl}
\multicolumn{2}{l}{\textbf{IMPERATIV}}\\
& \textbf{verschreibe (du)!}\\
& \textbf{verschreiben wir!}\\
& \textbf{verschreibt (ihr)!}\\
& \textbf{verschreiben Sie!}\\
\end{tabular}

\vspace{10pt}

\begin{tabular}{rl}
\multicolumn{2}{l}{\textbf{PERFEKT}}\\
& haben + \textbf{verschrieben}\\
\end{tabular}
\hfill
\begin{tabular}{rl}
\multicolumn{2}{l}{\textbf{FUTURE I}}\\
& werden + \textbf{verschreiben}\\
\end{tabular}
\hfill
\begin{tabular}{rl}
\multicolumn{2}{l}{\textbf{PLUSQUAMPERFEKT}}\\
& hatten + \textbf{verschrieben}\\
\end{tabular}

\vspace{10pt}

\begin{tabular}{lll}
\multicolumn{3}{l}{\textbf{MIT PRÄPOSITIONEN}}\\
& \textbf{---} & \\
\end{tabular}
\hfill
\begin{tabular}{lll}
\multicolumn{3}{l}{\textbf{TRENNBARE VERBEN}}\\
& \textbf{---} & \\
\end{tabular}
\vspace{-5pt}
\end{verbentry}

% verschwenden (to waste)
\begin{verbentry}{
  style = {default},
  toc = {verschwenden (to waste)},
  name = {verschwenden},
  english = {to waste},
  type = {regular},
  auxiliary = {haben}
}
 
    
     \vspace{5pt}

    \hrule
    
    \vspace{5pt}

  \begin{tabular}{rl}
     \multicolumn{2}{l}{\textbf{PRÄSENS} }\\
    ich & \textbf{verschwende}\\
    du & \textbf{verschwendest}\\
    er/sie/es & \textbf{verschwendet}\\
    wir & \textbf{verschwenden}\\
    ihr & \textbf{verschwendet}\\
    sie/Sie & \textbf{verschwenden}\\
  \end{tabular}
  \hfill
     \begin{tabular}{rl}
    \multicolumn{2}{l}{\textbf{PRÄTERITUM} }\\
    ich & \textbf{verschwendete}\\
    du & \textbf{verschwendetest}\\
    er/sie/es & \textbf{verschwendete}\\
    wir & \textbf{verschwendeten}\\
    ihr & \textbf{verschwendetet}\\
    sie/Sie & \textbf{verschwendeten}\\
  \end{tabular}
  \hfill
  \begin{tabular}{rl}
     \multicolumn{2}{l}{\textbf{IMPERATIV} }\\
   & \textbf{verschwende (du)!}\\
   & \textbf{verschwenden wir!}\\
   & \textbf{verschwendet (ihr)!}\\
   & \textbf{verschwenden Sie!}\\
   & \vspace{10pt}
  \end{tabular}

    \vspace{10pt}

 \begin{tabular}{rl}
     \multicolumn{2}{l}{\textbf{PERFEKT}}\\
   & haben + \textbf{verschwendet}\\
  \end{tabular}
\hfill
   \begin{tabular}{rl}
    \multicolumn{2}{l}{\textbf{FUTURE I} }\\
   & werde + \textbf{verschwenden}\\
  \end{tabular}
\hfill
   \begin{tabular}{rl}
      \multicolumn{2}{l}{\textbf{PLUSQUAMPERFEKT}}\\
   & hatten + \textbf{verschwendet}\\
  \end{tabular}

   \vspace{10pt}

   \begin{tabular}{lll}
      \multicolumn{3}{l}{\textbf{MIT PRÄPOSITIONEN}}\\
& \textbf{verschwenden an} + Akk & (to waste on)\\
& \textbf{verschwenden für} + Akk & (to waste for)\\
\end{tabular}
  \hfill
   \begin{tabular}{lll}
\multicolumn{3}{l}{\textbf{TRENNBARE VERBEN}}\\
& \textbf{---} & \\
\end{tabular}

  \vspace{-5pt}
\end{verbentry}

% verschwinden (to disappear)
\begin{verbentry}{
  style = {default},
  toc = {verschwinden (to disappear)},
  name = {verschwinden},
  english = {to disappear},
  type = {irregular},
  auxiliary = {sein}
}
 
    
     \vspace{5pt}

    \hrule
    
    \vspace{5pt}

  \begin{tabular}{rl}
     \multicolumn{2}{l}{\textbf{PRÄSENS} }\\
    ich & \textbf{verschwinde}\\
    du & \textbf{verschwindest}\\
    er/sie/es & \textbf{verschwindet}\\
    wir & \textbf{verschwinden}\\
    ihr & \textbf{verschwindet}\\
    sie/Sie & \textbf{verschwinden}\\
  \end{tabular}
  \hfill
     \begin{tabular}{rl}
    \multicolumn{2}{l}{\textbf{PRÄTERITUM} }\\
    ich & \textbf{verschwand}\\
    du & \textbf{verschwandest}\\
    er/sie/es & \textbf{verschwand}\\
    wir & \textbf{verschwanden}\\
    ihr & \textbf{verschwandet}\\
    sie/Sie & \textbf{verschwanden}\\
  \end{tabular}
  \hfill
  \begin{tabular}{rl}
     \multicolumn{2}{l}{\textbf{IMPERATIV} }\\
   & \textbf{verschwind(e) (du)!}\\
   & \textbf{verschwinden wir!}\\
   & \textbf{verschwindet (ihr)!}\\
   & \textbf{verschwinden Sie!}\\
   & \vspace{10pt}
  \end{tabular}

    \vspace{10pt}

 \begin{tabular}{rl}
     \multicolumn{2}{l}{\textbf{PERFEKT}}\\
   & sein + \textbf{verschwunden}\\
  \end{tabular}
\hfill
   \begin{tabular}{rl}
    \multicolumn{2}{l}{\textbf{FUTURE I} }\\
   & werde + \textbf{verschwinden}\\
  \end{tabular}
\hfill
   \begin{tabular}{rl}
      \multicolumn{2}{l}{\textbf{PLUSQUAMPERFEKT}}\\
   & waren + \textbf{verschwunden}\\
  \end{tabular}

   \vspace{10pt}

   \begin{tabular}{lll}
      \multicolumn{3}{l}{\textbf{MIT PRÄPOSITIONEN}}\\
& \textbf{verschwinden aus} + Dat & (to disappear from)\\
& \textbf{verschwinden in} + Dat & (to disappear in / into)\\
\end{tabular}
  \hfill
   \begin{tabular}{lll}
\multicolumn{3}{l}{\textbf{TRENNBARE VERBEN}}\\
& \textbf{---} & \\
\end{tabular}

  \vspace{-5pt}
\end{verbentry}

% verstehen (to understand)
\begin{verbentry}{
  style = {acc},
  toc = {verstehen (to understand)},
  name = {verstehen},
  english = {to understand},
  type = {irregular, + Akkusativ},
  auxiliary = {haben}
}
 
    
     \vspace{5pt}

    \hrule
    
    \vspace{5pt}

  \begin{tabular}{rl}
     \multicolumn{2}{l}{\textbf{PRÄSENS} }\\
    ich & \textbf{verstehe}\\
    du & \textbf{verstehst}\\
    er/sie/es & \textbf{versteht}\\
    wir & \textbf{verstehen}\\
    ihr & \textbf{versteht}\\
    sie/Sie & \textbf{verstehen}\\
  \end{tabular}
  \hfill
     \begin{tabular}{rl}
    \multicolumn{2}{l}{\textbf{PRÄTERITUM} }\\
    ich & \textbf{verstand}\\
    du & \textbf{verstandest}\\
    er/sie/es & \textbf{verstand}\\
    wir & \textbf{verstanden}\\
    ihr & \textbf{verstandet}\\
    sie/Sie & \textbf{verstanden}\\
  \end{tabular}
  \hfill
  \begin{tabular}{rl}
     \multicolumn{2}{l}{\textbf{IMPERATIV} }\\
   & \textbf{versteh(e) (du)!}\\
   & \textbf{verstehen wir!}\\
   & \textbf{versteht (ihr)!}\\
   & \textbf{verstehen Sie!}\\
   & \vspace{10pt}
  \end{tabular}

    \vspace{10pt}

 \begin{tabular}{rl}
     \multicolumn{2}{l}{\textbf{PERFEKT}}\\
  &  haben + \textbf{verstanden}\\
  \end{tabular}
\hfill
   \begin{tabular}{rl}
   
    \multicolumn{2}{l}{\textbf{FUTURE I} }\\
  &  werden + \textbf{verstehen}\\
 
  \end{tabular}
\hfill
   \begin{tabular}{rl}
      \multicolumn{2}{l}{\textbf{PLUSQUAMPERFEKT}}\\
  &  hatten + \textbf{verstanden}\\
  \end{tabular}
  
  
\vspace{10pt}

\begin{tabular}{lll}
\multicolumn{3}{l}{\textbf{MIT PRÄPOSITIONEN}}\\
& \textbf{---} & \\
\end{tabular}
\hfill
\begin{tabular}{lll}
\multicolumn{3}{l}{\textbf{TRENNBARE VERBEN}}\\
& \textbf{---} & \\
\end{tabular}
\vspace{-5pt}
\end{verbentry}

% versuchen (to try)
\begin{verbentry}{
  style = {acc},
  toc = {versuchen (to try)},
  name = {versuchen},
  english = {to try},
  type = {irregular, + Akkusativ},
  auxiliary = {haben}
}
 
    
     \vspace{5pt}

    \hrule
    
    \vspace{5pt}

  \begin{tabular}{rl}
     \multicolumn{2}{l}{\textbf{PRÄSENS} }\\
    ich & \textbf{versuche}\\
    du & \textbf{versuchst}\\
    er/sie/es & \textbf{versucht}\\
    wir & \textbf{versuchen}\\
    ihr & \textbf{versucht}\\
    sie/Sie & \textbf{versuchen}\\
  \end{tabular}
  \hfill
     \begin{tabular}{rl}
    \multicolumn{2}{l}{\textbf{PRÄTERITUM} }\\
    ich & \textbf{versuchte}\\
    du & \textbf{versuchtest}\\
    er/sie/es & \textbf{versuchte}\\
    wir & \textbf{versuchten}\\
    ihr & \textbf{versuchtet}\\
    sie/Sie & \textbf{versuchten}\\
  \end{tabular}
  \hfill
  \begin{tabular}{rl}
     \multicolumn{2}{l}{\textbf{IMPERATIV} }\\
   & \textbf{versuch(e) (du)!}\\
   & \textbf{versuchen wir!}\\
   & \textbf{versucht (ihr)!}\\
   & \textbf{versuchen Sie!}\\
   & \vspace{10pt}
  \end{tabular}

    \vspace{10pt}

 \begin{tabular}{rl}
     \multicolumn{2}{l}{\textbf{PERFEKT}}\\
  &  haben + \textbf{versucht}\\
  \end{tabular}
\hfill
   \begin{tabular}{rl}
   
    \multicolumn{2}{l}{\textbf{FUTURE I} }\\
  &  werden + \textbf{versuchen}\\
 
  \end{tabular}
\hfill
   \begin{tabular}{rl}
      \multicolumn{2}{l}{\textbf{PLUSQUAMPERFEKT}}\\
   & hatten + \textbf{versucht}\\
  \end{tabular}
  
  
\vspace{10pt}

\begin{tabular}{lll}
\multicolumn{3}{l}{\textbf{MIT PRÄPOSITIONEN}}\\
& \textbf{---} & \\
\end{tabular}
\hfill
\begin{tabular}{lll}
\multicolumn{3}{l}{\textbf{TRENNBARE VERBEN}}\\
& \textbf{---} & \\
\end{tabular}
\vspace{-5pt}
\end{verbentry}

% vertiefen (to deepen)
\begin{verbentry}{
  style = {acc},
  toc = {vertiefen (to deepen)},
  name = {vertiefen},
  english = {to deepen},
  type = {regular, + Akkusativ},
  auxiliary = {haben}
}


\vspace{5pt}
\hrule
\vspace{5pt}

\begin{tabular}{rl}
\multicolumn{2}{l}{\textbf{PRÄSENS}}\\
ich & \textbf{vertiefe}\\
du & \textbf{vertiefst}\\
er/sie/es & \textbf{vertieft}\\
wir & \textbf{vertiefen}\\
ihr & \textbf{vertieft}\\
sie/Sie & \textbf{vertiefen}\\
\end{tabular}
\hfill
\begin{tabular}{rl}
\multicolumn{2}{l}{\textbf{PRÄTERITUM}}\\
ich & \textbf{vertiefte}\\
du & \textbf{vertieftest}\\
er/sie/es & \textbf{vertiefte}\\
wir & \textbf{vertieften}\\
ihr & \textbf{vertieftet}\\
sie/Sie & \textbf{vertieften}\\
\end{tabular}
\hfill
\begin{tabular}{rl}
\multicolumn{2}{l}{\textbf{IMPERATIV}}\\
& \textbf{vertiefe (du)!}\\
& \textbf{vertiefen wir!}\\
& \textbf{vertieft (ihr)!}\\
& \textbf{vertiefen Sie!}\\
\end{tabular}

\vspace{10pt}

\begin{tabular}{rl}
\multicolumn{2}{l}{\textbf{PERFEKT}}\\
& haben + \textbf{vertieft}\\
\end{tabular}
\hfill
\begin{tabular}{rl}
\multicolumn{2}{l}{\textbf{FUTURE I}}\\
& werden + \textbf{vertiefen}\\
\end{tabular}
\hfill
\begin{tabular}{rl}
\multicolumn{2}{l}{\textbf{PLUSQUAMPERFEKT}}\\
& hatten + \textbf{vertieft}\\
\end{tabular}

\vspace{10pt}

\begin{tabular}{lll}
\multicolumn{3}{l}{\textbf{MIT PRÄPOSITIONEN}}\\
& \textbf{vertiefen in} + Akk & (to deepen into / immerse in)\\
\end{tabular}
\hfill
\begin{tabular}{lll}
\multicolumn{3}{l}{\textbf{TRENNBARE VERBEN}}\\
& \textbf{---} & \\
\end{tabular}

\vspace{-5pt}
\end{verbentry}

% vertrauen (to trust)
\begin{verbentry}{
  style = {default},
  toc = {vertrauen (to trust)},
  name = {vertrauen},
  english = {to trust},
  type = {regular},
  auxiliary = {haben}
}


\vspace{5pt}
\hrule
\vspace{5pt}

\begin{tabular}{rl}
\multicolumn{2}{l}{\textbf{PRÄSENS}}\\
ich & \textbf{vertraue}\\
du & \textbf{vertraust}\\
er/sie/es & \textbf{vertraut}\\
wir & \textbf{vertrauen}\\
ihr & \textbf{vertraut}\\
sie/Sie & \textbf{vertrauen}\\
\end{tabular}
\hfill
\begin{tabular}{rl}
\multicolumn{2}{l}{\textbf{PRÄTERITUM}}\\
ich & \textbf{vertraute}\\
du & \textbf{vertrautest}\\
er/sie/es & \textbf{vertraute}\\
wir & \textbf{vertrauten}\\
ihr & \textbf{vertrautet}\\
sie/Sie & \textbf{vertrauten}\\
\end{tabular}
\hfill
\begin{tabular}{rl}
\multicolumn{2}{l}{\textbf{IMPERATIV}}\\
& \textbf{vertraue (du)!}\\
& \textbf{vertrauen wir!}\\
& \textbf{vertraut (ihr)!}\\
& \textbf{vertrauen Sie!}\\
\end{tabular}

\vspace{10pt}

\begin{tabular}{rl}
\multicolumn{2}{l}{\textbf{PERFEKT}}\\
& haben + \textbf{vertraut}\\
\end{tabular}
\hfill
\begin{tabular}{rl}
\multicolumn{2}{l}{\textbf{FUTURE I}}\\
& werden + \textbf{vertrauen}\\
\end{tabular}
\hfill
\begin{tabular}{rl}
\multicolumn{2}{l}{\textbf{PLUSQUAMPERFEKT}}\\
& hatten + \textbf{vertraut}\\
\end{tabular}

\vspace{10pt}

\begin{tabular}{lll}
\multicolumn{3}{l}{\textbf{MIT PRÄPOSITIONEN}}\\
& \textbf{vertrauen auf} + Akk & (to rely on)\\
& \textbf{vertrauen in} + Akk & (to trust in)\\
& \textbf{vertrauen} + Dat & (to trust someone)\\
\end{tabular}
\hfill
\begin{tabular}{lll}
\multicolumn{3}{l}{\textbf{TRENNBARE VERBEN}}\\
\vspace{30pt}\\
\end{tabular}
\vspace{-5pt}
\end{verbentry}

% verzeihen (to forgive)
\begin{verbentry}{
  style = {dat},
  toc = {verzeihen (to forgive)},
  name = {verzeihen},
  english = {to forgive},
  type = {irregular, + Dativ},
  auxiliary = {haben}
}
 
    
     \vspace{5pt}

    \hrule
    
    \vspace{5pt}

  \begin{tabular}{rl}
     \multicolumn{2}{l}{\textbf{PRÄSENS} }\\
    ich & \textbf{verzeihe}\\
    du & \textbf{verzeihst}\\
    er/sie/es & \textbf{verzeiht}\\
    wir & \textbf{verzeihen}\\
    ihr & \textbf{verzeiht}\\
    sie/Sie & \textbf{verzeihen}\\
  \end{tabular}
  \hfill
     \begin{tabular}{rl}
    \multicolumn{2}{l}{\textbf{PRÄTERITUM} }\\
    ich & \textbf{verzieh}\\
    du & \textbf{verziehst}\\
    er/sie/es & \textbf{verzieh}\\
    wir & \textbf{verziehen}\\
    ihr & \textbf{verzieht}\\
    sie/Sie & \textbf{verziehen}\\
  \end{tabular}
  \hfill
  \begin{tabular}{rl}
     \multicolumn{2}{l}{\textbf{IMPERATIV} }\\
   & \textbf{verzeih(e) (du)!}\\
   & \textbf{verzeihen wir!}\\
   & \textbf{verzeiht (ihr)!}\\
   & \textbf{verzeihen Sie!}\\
   & \vspace{10pt}
  \end{tabular}

    \vspace{10pt}

 \begin{tabular}{rl}
     \multicolumn{2}{l}{\textbf{PERFEKT}}\\
   & haben + \textbf{verziehen}\\
  \end{tabular}
\hfill
   \begin{tabular}{rl}
   
    \multicolumn{2}{l}{\textbf{FUTURE I} }\\
   & werden + \textbf{verzeihen}\\
 
  \end{tabular}
\hfill
   \begin{tabular}{rl}
      \multicolumn{2}{l}{\textbf{PLUSQUAMPERFEKT}}\\
   & hatten + \textbf{verziehen}\\
  \end{tabular}

   \vspace{10pt}

   \begin{tabular}{lll}
      \multicolumn{3}{l}{\textbf{MIT PRÄPOSITIONEN}}\\
 & -- & --\\
\end{tabular}
  \hfill
   \begin{tabular}{lll}
      \multicolumn{3}{l}{\textbf{TRENNBARE VERBEN}}\\
 & -- & --\\
  \end{tabular}


  \vspace{-5pt}
\end{verbentry}

% verzichten (to do without / renounce)
\begin{verbentry}{
  style = {default},
  toc = {verzichten (to do without / renounce)},
  name = {verzichten},
  english = {to do without},
  type = {regular},
  auxiliary = {haben}
}


\vspace{5pt}
\hrule
\vspace{5pt}

\begin{tabular}{rl}
\multicolumn{2}{l}{\textbf{PRÄSENS}}\\
ich & \textbf{verzichte}\\
du & \textbf{verzichtest}\\
er/sie/es & \textbf{verzichtet}\\
wir & \textbf{verzichten}\\
ihr & \textbf{verzichtet}\\
sie/Sie & \textbf{verzichten}\\
\end{tabular}
\hfill
\begin{tabular}{rl}
\multicolumn{2}{l}{\textbf{PRÄTERITUM}}\\
ich & \textbf{verzichtete}\\
du & \textbf{verzichtetest}\\
er/sie/es & \textbf{verzichtete}\\
wir & \textbf{verzichteten}\\
ihr & \textbf{verzichtetet}\\
sie/Sie & \textbf{verzichteten}\\
\end{tabular}
\hfill
\begin{tabular}{rl}
\multicolumn{2}{l}{\textbf{IMPERATIV}}\\
& \textbf{verzichte (du)!}\\
& \textbf{verzichten wir!}\\
& \textbf{verzichtet (ihr)!}\\
& \textbf{verzichten Sie!}\\
\end{tabular}

\vspace{10pt}

\begin{tabular}{rl}
\multicolumn{2}{l}{\textbf{PERFEKT}}\\
& haben + \textbf{verzichtet}\\
\end{tabular}
\hfill
\begin{tabular}{rl}
\multicolumn{2}{l}{\textbf{FUTURE I}}\\
& werden + \textbf{verzichten}\\
\end{tabular}
\hfill
\begin{tabular}{rl}
\multicolumn{2}{l}{\textbf{PLUSQUAMPERFEKT}}\\
& hatten + \textbf{verzichtet}\\
\end{tabular}

\vspace{10pt}

\begin{tabular}{lll}
\multicolumn{3}{l}{\textbf{MIT PRÄPOSITIONEN}}\\
& \textbf{verzichten auf} + Akk & (to give up / do without)\\
\end{tabular}
\hfill
\begin{tabular}{lll}
\multicolumn{3}{l}{\textbf{TRENNBARE VERBEN}}\\
& \textbf{---} & \\
\end{tabular}
\vspace{-5pt}
\end{verbentry}

% verzieren (to decorate / adorn)
\begin{verbentry}{
  style = {acc},
  toc = {verzieren (to decorate / adorn)},
  name = {verzieren},
  english = {to decorate / adorn},
  type = {regular, + Akkusativ},
  auxiliary = {haben}
}


\vspace{5pt}
\hrule
\vspace{5pt}

\begin{tabular}{rl}
\multicolumn{2}{l}{\textbf{PRÄSENS}}\\
ich & \textbf{verziere}\\
du & \textbf{verzierst}\\
er/sie/es & \textbf{verziert}\\
wir & \textbf{verzieren}\\
ihr & \textbf{verziert}\\
sie/Sie & \textbf{verzieren}\\
\end{tabular}
\hfill
\begin{tabular}{rl}
\multicolumn{2}{l}{\textbf{PRÄTERITUM}}\\
ich & \textbf{verzierte}\\
du & \textbf{verziertest}\\
er/sie/es & \textbf{verzierte}\\
wir & \textbf{verzierten}\\
ihr & \textbf{verziertet}\\
sie/Sie & \textbf{verzierten}\\
\end{tabular}
\hfill
\begin{tabular}{rl}
\multicolumn{2}{l}{\textbf{IMPERATIV}}\\
& \textbf{verziere (du)!}\\
& \textbf{verzieren wir!}\\
& \textbf{verziert (ihr)!}\\
& \textbf{verzieren Sie!}\\
\end{tabular}

\vspace{10pt}

\begin{tabular}{rl}
\multicolumn{2}{l}{\textbf{PERFEKT}}\\
& haben + \textbf{verziert}\\
\end{tabular}
\hfill
\begin{tabular}{rl}
\multicolumn{2}{l}{\textbf{FUTURE I}}\\
& werden + \textbf{verzieren}\\
\end{tabular}
\hfill
\begin{tabular}{rl}
\multicolumn{2}{l}{\textbf{PLUSQUAMPERFEKT}}\\
& hatten + \textbf{verziert}\\
\end{tabular}

\vspace{10pt}

\begin{tabular}{lll}
\multicolumn{3}{l}{\textbf{MIT PRÄPOSITIONEN}}\\
& \textbf{verzieren mit} + Dat & (to decorate with)\\
\end{tabular}
\hfill
\begin{tabular}{lll}
\multicolumn{3}{l}{\textbf{TRENNBARE VERBEN}}\\
& \textbf{---} & \\
\end{tabular}

\vspace{-5pt}
\end{verbentry}
